\documentclass[11pt]{article}
\usepackage[T1]{fontenc}
\usepackage{textcomp}
\usepackage[margin=0.75in]{geometry}
\usepackage{amsmath,amssymb,amsthm}
\usepackage{array,booktabs}
\usepackage{hyperref}
\setlength{\parindent}{0pt}
\newtheorem{theorem}{Theorem}
\theoremstyle{definition}
\newtheorem{definition}{Definition}
\title{Flow Proof Artifact: prop-03-cut-equal-segment}
\author{Generated by \texttt{flow doc proof}}
\date{Flow Proof Book}
\begin{document}
\maketitle
Given two unequal straight lines, to cut off from the greater a part equal to the less.\\[0.5em]
\textbf{Source.} euclid — Elements, Book I, Proposition 3\\[0.5em]
\bigskip
\noindent{\large \textbf{Derived fact 1.} \textit{Given two unequal straight lines, to cut off from the greater a part equal to the less} \hfill the Euclidean plane \cdot Euclid Book I \cdot Proposition 3: cut a segment equal to a smaller segment}\par
\smallskip\noindent\textit{Coordinate: the Euclidean plane \cdot Euclid Book I \cdot Proposition 3: cut a segment equal to a smaller segment}\par
\smallskip\noindent\textit{Source: euclid — Elements, Book I, Proposition 3}\par
\smallskip\noindent\textit{Built on: proposition 1: equilateral triangle on a segment, for Book I of the Elements in on the Euclidean plane, proposition 2: a segment equal to a given segment, for Book I of the Elements in on the Euclidean plane}\par
\medskip
\noindent\textbf{Goal.} Given two unequal straight lines, to cut off from the greater a part equal to the less\par
\noindent$\displaystyle \text{cut from AE equals CD}$\par
\medskip
\noindent
\begin{tabular}{@{} >{\raggedright\arraybackslash}p{0.06\textwidth} >{\raggedright\arraybackslash}p{0.47\textwidth} >{\raggedleft\arraybackslash}p{0.41\textwidth} @{}}
\textbf{\#} & \textbf{Proof} & \textbf{Mathematics} \\
\midrule
\textbf{1.} & We prove that proposition 3: cut a segment equal to a smaller segment for Book I of the Elements in the Euclidean plane. & \\[0.45em]
\textbf{2.} & Let equilateral\_AEF = triangle\_on\_AE\_by\_Proposition\_1. & \\[0.45em]
\textbf{3.} & Let circle\_center\_A\_radius\_CD = circle\_at\_A\_with\_radius\_CD. & \\[0.45em]
\textbf{4.} & Let F = point\_where\_circle\_meets\_line\_from\_A\_through\_E. & \\[0.45em]
\textbf{5.} & We invoke the axiom governing Book I of the Elements in the Euclidean plane: proposition 1: equilateral triangle on a segment, for Book I of the Elements in on the Euclidean plane. & \\[0.45em]
\textbf{6.} & We invoke the axiom governing Book I of the Elements in the Euclidean plane: proposition 2: a segment equal to a given segment, for Book I of the Elements in on the Euclidean plane. & \\[0.45em]
\textbf{7.} & From step 2, step 3, step 4, step 5, and step 6, we can deduce that segment AF equals segment CD. & $\displaystyle \text{segment AF} = \text{segment CD}$ \\[0.45em]
\textbf{8.} & We can deduce that cut from AE equals CD. Hence proven. & $\displaystyle \text{cut from AE equals CD}$ \\[0.45em]
\bottomrule
\end{tabular}
\label{thm:Geometry-Euclid Book I-proposition_3_cut_a_segment_equal_to_a_smaller_segment}
\par\medskip\hrule
\end{document}
